% Part: second-order-logic
% Chapter: syntax-and-semantics
% Section: language-of-sol

\documentclass[../../../include/open-logic-section]{subfiles}

\begin{document}

\olfileid{sol}{syn}{lan}

\olsection{द्वितीय-क्रम तर्कशास्त्राची भाषा}

प्रथम-क्रम तर्कशास्त्राप्रमाणेच द्वितीय-क्रम तर्कशास्त्रातील
अभिव्यक्ती मूलभूत शब्दसंग्रहापासून बनतात. त्यात \emph{!!{variable}s},
\emph{!!{constant}s}, \emph{!!{predicate}s} आणि कधी कधी
\emph{!!{function}s} असतात. यांच्याबरोबर तार्किक संकारके,
संख्यक आणि कंस व स्वल्पविरामांसारखी विरामचिन्हे वापरून
\emph{पदे} आणि \emph{!!{formula}s} तयार होतात. फरक असा की
वस्तूंसाठीच्या चलांव्यतिरिक्त द्वितीय-क्रम तर्कशास्त्रात संबंध
आणि फलने यांसाठीही चले असतात आणि त्यांच्यावर संख्यापन
करता येते. म्हणून द्वितीय-क्रम तर्कशास्त्राची तार्किक चिन्हे
म्हणजे प्रथम-क्रम तर्कशास्त्राची चिन्हे, तसेच पुढील चिन्हे:

\begin{enumerate}
\item प्रत्येक स्थानसंख्या~$n$ साठी द्वितीय-क्रम संबंध !!{variable}s चा
  !!{denumerable}s संच: $\Obj V_0^n$, $\Obj V_1^n$, $\Obj V_2^n$, \dots
\item द्वितीय-क्रम फलन !!{variable}s चा !!{denumerable}s संच:
  $\Obj u_0^n$, $\Obj u_1^n$, $\Obj u_2^n$, \dots
\end{enumerate}

प्रथम-क्रम तर्कशास्त्रात !!{identity}~$\eq$ सहसा समाविष्ट
असते. प्रथम-क्रम तर्कशास्त्रात $\Lang{L}$ या भाषेतील
अतार्किक चिन्हे आपल्याला काही अर्थपूर्ण गोष्ट व्यक्त करता
यावी यासाठी अत्यावश्यक असतात. अतार्किक चिन्हे नसलेली
!!{sentences} नक्कीच आहेत, पण केवळ~$\eq$ वापरून
फारसे अर्थपूर्ण काही सांगता येत नाही. द्वितीय-क्रम
तर्कशास्त्रात मात्र संबंधचल आणि फलनचल अमर्याद प्रमाणात
उपलब्ध असल्याने, अतार्किक चिन्हांचा वेगळा साठा नसतानाही
प्रथम-क्रम भाषेत सांगता येणारी कोणतीही गोष्ट सांगता येते.

\end{document}
